%------------------------------------------------------------------------------------
%  SECTION 6.4 : LOGO DYNAMIQUE SCRLAYER ET DOCUMENTS PDF EXTERNES
%------------------------------------------------------------------------------------
\section{Logo en Calque Dynamique et Insertion de Documents Externes}
\label{sec:logo_documents_externes}

Deux défis récurrents dans la formalisation d'un livrable officiel concernent la présence d'un filigrane institutionnel discret et la fusion de pièces jointes administratives ou techniques fournies par des tiers.

\subsection{Le filigrane dynamique en arrière-plan (\textsf{scrlayer})}
\label{subsec:calque_logo}

L'insertion manuelle d'un logo dans les en-têtes (\texttt{\textbackslash ihead} ou \texttt{\textbackslash ohead}) modifie l'encombrement vertical (\texttt{headheight}) et crée des décalages de texte selon que la page est paire ou impaire.

\texttt{ClasseGuillaume} s'affranchit de cette contrainte en utilisant le concept de **couche d'arrière-plan** (\emph{background layer}) du moteur \textsf{scrlayer} :

\begin{verbatim}
\DeclareNewLayer[
  background,
  topmargin,
  align = tr,
  hoffset = \paperwidth - 1.8cm, % Calibré dynamiquement sur le bord droit
  voffset = 1.2cm,
  width = 3cm,
  height = 3cm,
  contents = {%
    \IfFileExists{images/logo_test.png}{%
      \includegraphics[width=2.2cm]{images/logo_test.png}%
    }{}%
  }
]{logo.layer}
\end{verbatim}

\begin{boiteSynthese}{Commandes de pilotage du logo institutionnel}
  Le calque flotte au-dessus de la mise en page sans jamais interférer avec le texte du corps. Deux commandes permettent d'en moduler l'affichage selon les chapitres :
  \begin{itemize}
    \item \texttt{\textbackslash masquerLogo} : Désactive instantanément le calque (utilisé dans le Frontmatter pour préserver la virginité de la page de titre et du résumé).
    \item \texttt{\textbackslash afficherLogo} : Réactive le calque en tête du corps de texte (\texttt{\textbackslash mainmatter}) ou dans les sections spécialisées.
  \end{itemize}
\end{boiteSynthese}

\subsection{Intégration transparente de documents PDF externes avec \textsf{pdfpages}}
\label{subsec:pdfpages_externes}

Dans une thèse ou un rapport d'ingénierie, il est fréquent de devoir intégrer des documents tiers au format PDF :
attestations administratives, comptes-rendus de laboratoires accrédités, fiches techniques de constructeurs, ou articles scientifiques complets déjà publiés.
Réinsérer ces documents sous forme d'images capturées dégrade fortement la résolution et perd la numérotation séquentielle des pages.

La classe intègre le paquet de référence \textsf{pdfpages} :
\begin{verbatim}
% Exemple d'inclusion d'un document multipage en annexe :
\includepdf[
  pages = {1-3},                  % Pages à importer
  scale = 0.92,                   % Échelle réservant la marge KOMA
  pagecommand = {\thispagestyle{scrheadings}} % En-têtes conservés
]{documents/rapport_mesures.pdf}
\end{verbatim}
Cette méthode vectorielle conserve la netteté parfaite des textes et graphiques d'origine tout en assurant une continuité sans faille de la pagination générale de votre mémoire.
Une démonstration concrète est proposée en Annexe~\ref{annexe:specimen_externe}.

